In this paper, we aimed to thoroughly investigate the current AI Scientist capabilities and the associated risks.
To this end, we developed Jr. AI Scientist, an AI Scientist specialized in extending a given baseline paper.
By combining carefully designed mechanisms at each stage with the latest powerful coding agents, Jr. AI Scientist is capable of autonomously generating research papers of higher quality than those produced by existing systems.
This provides valuable insights into the capability of our Jr. AI Scientist.
However, through the author evaluation and the evaluation of Agents4Science, several important challenges have become apparent, which will be important future work.
Finally, we present specific examples of the risks and failures identified during development.
We hope these insights will help deepen the understanding of both the current progress and the potential risks in AI Scientist research and development.

\section*{Limitations and Future Work}
\textbf{Use of Multiple Agents.}
We built our AI Scientist on top of a single coding agent (\ie~Claude Code) and a single family of large language models. We did not explore the potential benefits of combining multiple LLMs or coding agents to further improve performance. Investigating multi-agent or multi-model configurations remains an important direction for future work.

\textbf{Fragility of Novelty Verification via Semantic Scholar.}
Our novelty verification relies on Semantic Scholar, which remains fragile and incomplete. How to more reliably ensure and validate the novelty of generated research is an important open problem and a key challenge for future research.

\textbf{Addressing Observed Risks and Identified Limitations.}
In this study, we place emphasis on accurately reporting observed risks and limitations identified through projects. Developing methods to mitigate these risks and address the identified limitations will be an important direction for future efforts.


\section*{Broader Impact Statement}
Through this work, we developed a Jr. AI Scientist. However, we do not recommend the use of Jr. AI Scientist or similar systems in the preparation of actual conference submission papers. The primary objective of this paper is to help the community gain a clearer and more comprehensive understanding of the current status of AI Scientists. As discussed throughout the paper, AI Scientists entail various risks. Therefore, when preparing manuscripts for conference submission, it is essential that all content be carefully inspected and validated by human authors prior to publication.

Furthermore, we argue that researchers working on AI Scientists bear a responsibility to report risks, including concrete failure cases and negative outcomes. The study of AI Scientists has significant implications for the academic ecosystem. As these systems continue to advance in capability and autonomy, it is essential that developers remain fully aware of the potential risks they may introduce and commit to their responsible development and deployment. Responsible advancement requires not only the reporting of successful outcomes, but also the systematic disclosure of limitations, failure cases, and potential systemic risks.

This paper adheres to such a responsible research principle. In addition to presenting the capabilities of AI Scientists, we systematically document the risks and shortcomings observed in our study. Through this transparent reporting, we hope to help the community gain a clearer and more comprehensive understanding of the current status of AI Scientists.